% PM 트랙 Day 3 슬라이드 본문 — 손으로 쓴 정본(한 프레임 = 한 쪽). Day2_슬라이드_body.tex 와 같은 형식.
% 래퍼: Day3_슬라이드.tex(슬라이드판) / Day3_슬라이드_노트.tex(노트판, 우측에 \note).
% 화면에는 결론만, 설명·근거·예외는 각 프레임의 \note{}에. Day3_슬라이드.md 는 이 파일의 텍스트 미러다.
% 그림은 ../그림/PM_Day3/*.pdf (벡터). 그림 번호 = 파일명 번호(그림 2-3 = fig_2_3_*). 모든 숫자는 Day3 워크북 완성 예시본(정본)에서.
% 데이터 일자 2026-09-30. BAC 145.80 / 집행 한도 156.00 / PV 77.83 / EV 73.43 / AC 71.13 / SPI 0.943 / SPI(t) 0.920 / CPI 1.032(조정 0.968) / EAC4 150.86.


\newcolumntype{L}[1]{>{\raggedright\arraybackslash}p{#1}}
\newcommand{\ra}{$\rightarrow$}
\newcommand{\til}{\textasciitilde}
\newcommand{\keymsg}[1]{\par\vfill{\color{mDarkTeal}\bfseries #1}\par}
\newcommand{\figcap}[1]{\par\smallskip{\footnotesize\color{black!60}#1}\par}
\newcommand{\fig}[2][0.66]{\centering\includegraphics[height=#1\textheight,width=\textwidth,keepaspectratio]{#2}\par}
\newcommand{\subhead}[1]{{\color{mDarkTeal}\bfseries #1}\par}
\newenvironment{tbl}[2][\small]{\begin{center}\usebeamercolor[fg]{normal text}\setlength{\tabcolsep}{4pt}\renewcommand{\arraystretch}{1.15}#1\begin{tabular}{#2}\toprule}{\bottomrule\end{tabular}\end{center}}

% =====================================================================
% 도입

% =====================================================================
% 도입
% =====================================================================
\begin{frame}{PM 트랙 Day 3 — 어제 세운 자(尺)에 오늘 실적을 댄다}
\begin{itemize}
\item 프로젝트 매니지먼트 4일 특강 · 3일차 · 실행과 통제(Executing · Monitoring \& Controlling)
\item 사례는 같다 — 온담F\&B홀딩스 ONE ONDAM \textbf{P1} · 승인 168.0 / 집행 한도 156.0 / PMB 145.8억 · FP 12,400 고정가
\item 시점 고정 — \textbf{데이터 일자 2026-09-30}. 착수 9개월, S19 진행 중(R2 완료 직후), 규제 ① 검증 통과 직후, G2(11-27) 앞
\item 어제 등록부의 리스크가 실제로 일어났다 — R-05 부분 실현(구현 3주 지연), R-04 대응(라이선스 50/50), R-03 트리거 발동(폴백 확정)
\item Day2 다섯 기준선(⑥\til{}⑩)이 오늘 다섯 통제 산출물(⑪\til{}⑮)의 \textbf{자}다
\end{itemize}
\keymsg{오늘의 한 문장 — 통제는 편차를 재는 일이 아니라, 편차에 대해 누가 무엇을 결정해야 하는지를 적는 일이다}
\note{표지. 어제 워크북 완성 예시본 5종 — 특히 §3.5 CPM 계산표·스프린트 캘린더, §4.5 월별 PV 표(2026-09 누적 77.83), §5.5 등록부(잔여 EMV 6.78)와 격자 7축 — 를 읽어 왔는지 확인한다. 읽지 않은 수강생은 워크북 §0.3 한 페이지 요약(2026-09-30까지 일어난 일 아홉 줄)과 §0.4 "Day2에서 가져오는 것" 표를 4교시 전까지 읽도록 안내한다. Day1이 "결정되지 않은 것을 결정하는 문서", Day2가 "결정된 것을 분해하고 숫자로 만드는 문서"였다면, Day3는 "기준선과 실적의 차이를 재고 그 차이에 대해 누가 언제까지 무엇을 결정해야 하는지를 적는 문서"다. 그래서 오늘의 모든 산출물은 결정 요청 칸으로 끝난다. 시나리오는 [추가 설정]이며 숫자는 Day2 기준선과 정합하게 정했다 — 워크북과 교재가 다르면 워크북이 정본. 예상 소요 3분.}
\end{frame}

\begin{frame}{Day2 기준선 5종 \ra{} Day3 통제 산출물 5종}
\fig[0.62]{fig_0_1_day2_to_day3}
\figcap{그림 0-1. 성과보고(⑫)가 중심이고 나머지 넷이 그 각주다. 되돌아가는 파선은 승인된 변경(BL-02)으로만 기준선이 움직였음을 보인다 — 지연(A10 $-$10)과 초과(밀도 0.132)는 편차로 남아 있다}
\keymsg{Day1 §0 마스터 표의 Day3 3종이 Day2 §7의 요구(품질 통제)를 받아 5종으로 확정되었다 — 4일 20종}
\note{교재 서문의 표와 그림 0-1. ⑪ 변경요청서·영향 분석·CCB 의결서는 ⑥ WBS 사전 카드(추정·의존·인수 기준)·⑦ RTM 변경 이력 열·⑩ 격자(±35건·FP ±372·원가 10.2·리스크 8억)와 Day2 §6의 헌장 v1.1 수정 요청 10건 중 1·2·4번을 받는다. ⑫ 성과보고는 ⑨ 월별 PV 표 17개월(BAC = PMB 145.8)과 ⑧ CPM 계산표·스프린트 캘린더를. ⑬ 이슈 로그·리스크 등록부 갱신은 ⑩ 등록부(잔여 6.78·트리거·2차 리스크 예비 번호 R-14\til{}R-18)와 컨틴전시 배정표(4.60 + 5.60)를. ⑭ 품질 통제 기록은 ⑦ 시험·검수 열과 ⑩ 품질 허용범위 두 계층(작업패키지 0.13 / 프로젝트 0.4 + 심각도 1·2 미해결 0)을. ⑮ 조달·계약 변경 기록은 ⑨ 단가 55만 원/FP·⑥ L2 FP 배분·⑧ 여유 소유권 문안을. 세 가지를 미리 말한다 — 기준선 값을 바꾸지 말고 인용하라(틀렸음을 발견하면 "측정 규칙 차이"로 분리해 적고 CR-10으로 CCB에) / 오늘의 시나리오는 나쁜 소식을 담고 있다(좋은 소식만 담은 사례는 연습이 안 된다) / 모든 문서는 결정 요청으로 끝난다. 예상 소요 4분.}
\end{frame}

\begin{frame}{오늘의 흐름}
\begin{columns}[T]
\begin{column}{0.47\textwidth}
\subhead{오전 — 이론 (제1\til{}3장)}
\begin{itemize}
\item \textbf{1교시} 90분 · 통제론 · 성과 측정 EVM — PMB·BAC·Rules of Credit·편차와 지수 — 제1장, 제2장 2.1\til{}2.3
\item \textbf{2교시} 90분 · EAC 4공식 · Earned Schedule · 착시 5개 · 흐름 지표 · 임계경로 $-$10과 EX-01 — 제2장 2.3\til{}2.9
\item \textbf{3교시} 90분 · 변경 관리 — 6축 · 처리 경로 · BL-02 vs 재기준선 · 변경 대가 — 제3장
\end{itemize}
\end{column}
\begin{column}{0.51\textwidth}
\subhead{오후 — 실습과 통제의 나머지}
\begin{itemize}
\item \textbf{4교시} 90분 · \textbf{실습 ①} 변경요청서·영향 분석·CCB + 성과보고 EVM 손계산 — 워크북 §1\til{}2
\item \textbf{5교시} 75분 · 품질 · 이슈·리스크 · 조달 통제 — 결함 곡선·감리 귀속·EMV 재평가·검수 보류 회신·귀책 분해 — 제4\til{}5장
\item \textbf{6교시} 75분 · \textbf{실습 ②} 이슈 로그·품질 기록·감리 대응·검수 보류 회신 — 워크북 §3\til{}5
\item \textbf{마무리} 30분 · 학술 배경 · 읽을 자료 · Day4 예고 — 제6장
\end{itemize}
\end{column}
\end{columns}
\keymsg{오전 세 교시가 성과보고와 변경요청서의 이론이고, 그 둘을 4교시에 손으로 쓴다 — 나머지 셋은 오후에}
\note{Day2와 같은 골격 — 오전 이론 셋, 점심은 3교시 뒤, 오후는 실습 ① \ra{} 이론(품질·이슈·조달) \ra{} 실습 ②. 1교시가 제2장 전반(2.1\til{}2.3 편차·지수까지)을 포함하는 이유는 EVM의 정의와 Rules of Credit이 오전 내내 쓰이기 때문이고, 2교시는 EAC 4공식부터 시작한다. 4교시 실습 ①이 ⑪·⑫ 둘을 맡는 것은 두 문서가 서로의 입력이기 때문 — 변경 4건의 FP 델타(+185)가 EAC4의 ETC 행에 들어가고, 성과보고의 "미반영 변경 CR-03·07"이 변경 로그를 가리킨다. 5교시 75분에 제4장(품질·이슈·리스크·예외 보고)과 제5장(보고 체계·mum effect·조달·계약)을 다 다루므로 속도가 빠르다 — 6교시 실습 ②가 그 내용을 손으로 다시 밟는다. 예상 소요 3분.}
\end{frame}

\begin{frame}{학습 목표 — 오늘이 끝나면}
\begin{columns}[T]
\begin{column}{0.5\textwidth}
\subhead{오전 이론}
\begin{enumerate}
\item 통제 사이클 다섯 단계와 조치의 세 층(PM 전결 / 예외 보고 / 변경요청)을 \textbf{허용범위 격자}로 가른다
\item BAC가 145.8인 이유를 ANSI/PMI 19-006 용어로 \textbf{설명한다}
\item EAC 4공식을 전제와 함께 계산하고, SPI가 오르는데 예외 보고를 내는 이유를 \textbf{SPI(t)와 CPM}으로 말한다
\item 착시 다섯에 이름을 붙이고 \textbf{병기}한다 — 대체하지 않는다
\item 변경을 6축으로 재고 처리 경로 규칙 ②(헌장 문언)로 CCB 여부를 판정한다
\end{enumerate}
\end{column}
\begin{column}{0.48\textwidth}
\subhead{오후}
\begin{enumerate}\setcounter{enumi}{5}
\item 결함 곡선의 평탄을 시험 노력으로 판독하고, 감리 지적의 \textbf{귀속}에 문서로 이의한다
\item 잔여 EMV 변동을 분해하고 컨틴전시가 \textbf{배정으로} 소진됨을 읽는다
\item 구두 검수 보류에 서면으로 답하고 지연 귀책을 \textbf{지금} 분해한다
\item 다섯 산출물을 온담 P1으로 \textbf{직접 작성}하고 서로 \textbf{설명}하게 한다
\end{enumerate}
\end{column}
\end{columns}
\keymsg{정직하게 적을 셋 — SPI 0.943이 오르는데 EX-01을 낸 이유 / CPI 1.032가 좋은 소식이 아닌 이유 / 소진 13.8\%가 안심이 아닌 이유}
\note{아홉 목표 중 1\til{}5는 오전(1\til{}3교시), 6\til{}8은 5교시, 9는 두 실습. 오늘의 결론을 미리 말해 둔다(교재 "Day3을 마치며") — 실행 단계에서 정직하게 적어야 할 세 가지. SPI가 0.913 \ra{} 0.943으로 오르는데 예외 보고를 낸 이유는 SPI(t)가 0.920으로 정체이고 임계경로 A10이 $-$10영업일이며 컷오버 창은 허용범위 0이라는 사실. CPI 1.032가 좋은 소식이 아닌 이유는 라이선스 4.70의 인식 기준 차이를 걷어 내면 0.968이고 EAC3$'$ 155.04는 집행 한도까지 0.96 남았다는 사실. 컨틴전시 소진 13.8\%가 안심의 근거가 아닌 이유는 미배정 풀이 스프린트 4회분에서 1.75회분이 되었고 관리예비비 대상 2.20의 원인이 전부 프로그램·스폰서 계층이라는 사실. 이 셋은 나쁜 소식이 아니라 통제가 한 일이며, 스폰서에게 결정을 요청하는 근거다. 그리고 재기준선은 0회다. 예상 소요 3분.}
\end{frame}

% =====================================================================
\section{1교시 (90분) — 통제론 · 성과 측정 EVM: PMB · BAC · Rules of Credit · 편차와 지수}
% =====================================================================

\begin{frame}{이 교시에서 배우는 것}
\begin{itemize}
\item 통제 사이클 다섯 단계 — 측정 \ra{} 분석 \ra{} 예측 \ra{} 조치 \ra{} 기준선 갱신, 그리고 조치의 세 층 [§1.1]
\item 통제의 단위 — 예측형의 게이트, 적응형의 반복, 하이브리드의 세 층(스프린트 / 월 / 게이트)과 층 사이의 연결 규칙 [§1.2]
\item 표준 — PMBOK 8판 두 Focus Area · PRINCE2 7 세 프로세스와 예외에 의한 관리 · ISO 21502 §7 · 한국 SI 네 주기 [§1.3\til{}1.4]
\item 가장 비싼 오류 셋 — 90\% 증후군 · 기준선 없는 보고 · 계획을 고쳐 편차를 없애기 [§1.5]
\item EVM의 전제 PMB와 BAC(145.8 vs 156.0) · Rules of Credit 여섯 기법 · SV·CV·SPI·CPI 읽기 규칙 셋 [§2.1\til{}2.3]
\end{itemize}
\keymsg{자가 있어야 편차가 있고, 편차가 있어야 통제가 있다 — 그러나 편차를 재는 것만으로는 아무 일도 일어나지 않는다}
\note{제1장 전체와 제2장 2.1\til{}2.3 편차·지수까지를 90분에 다룬다. 시간 배분은 §1.1 12분, §1.2 10분, §1.3\til{}1.4 12분, §1.5 8분, §2.1 15분, §2.2 15분, §2.3 편차·지수 10분, 정리 5분(EAC 4공식은 2교시 첫 주제). §2.1 BAC 논쟁과 §2.2 Rules of Credit은 4교시 실습 ① 성과보고 손계산의 직접 입력이므로 여기서 숫자를 정확히 남긴다 — PV 77.83 / EV 73.43 / AC 71.13. 수강생이 PMP 소지자라면 ANSI/PMI 19-006-2019가 2011년 Practice Standard 2판을 대체한 현행 정본임을 판본명으로 정확히 부르는 것이 신뢰의 첫 단계. 예상 소요 2분.}
\end{frame}

\begin{frame}{통제 사이클 — 측정 \ra{} 분석 \ra{} 예측 \ra{} 조치 \ra{} 기준선 갱신}
\fig[0.66]{fig_1_1_control_cycle}
\figcap{그림 1-1. 다섯 단계와 온담 2026-09-30의 실제 값. 조치의 세 층(PM 전결 / 예외 보고 / 변경요청)을 가르는 것이 Day1의 허용범위 격자, 기준선 갱신은 승인된 변경으로만. 바깥 동심원 = 통제의 주기}
\keymsg{"SV $-$4.40억"은 사실이지 결정이 아니다 — 통제는 편차에 대해 무엇을 할지를 정하고 실행하는 일이다}
\note{§1.1 개념. 측정 — 데이터 일자를 고정하고 기준선과 같은 규칙으로 잰다(인수 FP 7,020 · AC 71.13 · A10 착수 09-28 · 결함 1,140/1,024). 분석 — 크기가 아니라 원인: SV $-$4.40 중 $-$1.96은 측정 시차이고 $-$2.44가 실질 지연이며 그것은 SI 열($-$5.31)에 몰려 있고 시험·감리(+0.70)가 일부 가린다. 예측 — 전제를 적는다: EAC1 141.23은 "지금 효율이 끝까지", EAC4 150.86은 항목별 잔여 근거의 합. 조치 — 세 층: 허용범위 안 PM 전결(결함 해소 전담 2명), 초과 예상 시 예외 보고(EX-01 컷오버 창 이탈), 기준선 변경은 CR \ra{} CCB. 기준선 갱신 — CCB 제10차(10-15)가 CR-08\til{}11을 의결하고 BL-02(10-23)가 발행되면 FP 12,400 \ra{} 12,585, PMB 145.80 \ra{} 146.82; 그러나 A10 $-$10은 기준선에 들어가지 않는다. 주기의 길이(월 성과보고·격주 변경협의체·2주 리뷰·주간 PM 회의)가 통제의 해상도. 예상 소요 6분.}
\end{frame}

\begin{frame}{온담으로 읽기 — 사이클은 아홉 번 돌았고, 두 번 실패하고 한 번 성공했다}
\begin{itemize}
\item \textbf{실패 1 — 4월 제4호}: SPI 0.784를 "주의"로 적고 끝. 열별 분해가 없어 상용SW $-$6.0(조달 이연)과 SI 실질 0이 한 숫자로 — 측정은 되었으나 분석이 없었고, 그래서 조치도 없었다
\item \textbf{실패 2 — 5월 제5호}: CV가 $-$0.16 \ra{} +2.70. "원가 절감"으로 읽혔다. 실제는 03-13 라이선스 50/50 계약 변경의 흔적 — EV 100\% 인식, AC 50\%만 발생. 8월에야 발견
\item \textbf{성공 — 9월 제9호}: S18 3주 연장(09-04) \ra{} A10 착수 09-28 \ra{} TF $-$10 \ra{} as-is MG3 2027-03-05 \ra{} 컷오버 1차 창 이탈 예측 \ra{} \textbf{EX-01 10-02 제출}(인지 09-28 + 4영업일, 규정 5영업일 안)
\item 두 실패와 한 성공의 차이는 도구가 아니라 \textbf{양식}이다 — 4월에 열별 표가, 5월에 조정 행이 있었다면
\item 그래서 제9호는 조기 신호 5개를 사후에 적어 다음 주기의 양식을 바꿨다 — 조정 행 · 시험 노력 열 · "임계경로 TF 영향" 칸
\end{itemize}
\keymsg{통제 사이클의 실무적 정체는 보고서 양식이다 — 양식에 없는 칸은 관측되지 않는다}
\note{§1.1 "온담의 사례로 읽기". 4월의 실패에서 R-04 이슈 전환(IS-01)은 2월에 이미 있었지만 그것이 4월 SPI의 원인이라는 연결이 보고서에 없었다 — 분석 열이 없으면 이슈 로그와 성과보고는 서로를 설명하지 못한다. 5월의 실패는 측정 규칙(AC 인식 기준)이 실행 중에 바뀌었는데 성과보고 양식에 "조정 행"이 없었던 문제이며, 이 착시는 나중에 2027-01 자금 봉우리에 얹혀 R-17(2027 자금 집중)의 일부가 된다 — 착시는 지표만 왜곡하는 것이 아니라 리스크를 만든다. 9월의 성공은 스프린트 캘린더와 CPM 계산표가 같은 날짜를 가리켰기 때문에 가능했다(Day2 3.1절의 세 층). 다음 프레임에서 그 세 층을 본다. 예상 소요 5분.}
\end{frame}

\begin{frame}{통제의 단위 — 예측형의 게이트, 적응형의 반복, 하이브리드의 세 층}
\fig[0.6]{fig_1_2_three_control_layers}
\figcap{그림 1-2. 같은 사건(S18 3주 연장)이 세 층을 관통하며 다른 해상도로 읽힌다 — 스프린트 층 "이월 420 FP" / 월 층 "A10 TF $-$10, SPI(t) 0.920" / 게이트 층 "컷오버 1차 창 이탈 예상 \ra{} 예외 보고"}
\keymsg{층 사이의 연결 규칙이 없으면 아래층의 작은 결정이 위층의 큰 사건이 될 때까지 아무도 보지 못한다}
\note{§1.2 개념. 그림은 왼쪽에서 오른쪽으로 읽는다 — 맨 아래 스프린트 층(2주 × 36회: 스프린트 리뷰에서 인수 FP 확인, throughput·WIP·work item age, 결정은 이재현 PM과 P1 PMO), 가운데 월 층(성과보고 EVM·ES·CPM TF, 변경협의체 격주·CCB 월, 등록부 갱신, 요구사항 동결 매월 마지막 수요일 — 결정은 P1 PM 전결과 CCB), 맨 위 게이트 층(G0\til{}G5: 재기준선 여부, 컷오버 go/no-go, 헌장 확정 — 스폰서·프로그램 이사회). 예측형의 단위는 게이트와 보고 주기이고 EVM은 그 자연스러운 도구(PMB가 있어야 PV가 있다). 적응형의 단위는 반복과 리뷰이며 "기준선 대비 편차" 개념이 약하고, EVM을 그대로 쓰면 BAC가 흔들리고 스토리포인트는 EV 단위가 못 된다(2.7절). 하이브리드는 둘을 층으로 겹친다. 온담 사례 — 05-12 POS 화면 정의서 승인 대기 시작(R-05 실현): 스프린트 층은 S10을 3주로 늘려 흡수, 월 층은 SI SV $-$2.49를 "약간의 지연"으로 넘겼고 대체 결정권자 지정은 06-03(3주 늦음), 변경협의체는 07-15에 12영업일 중 7영업일 발주사 귀책·0.42억 판정, 게이트 층은 G2 앞 R-05 재평가로만 본다. 예상 소요 6분.}
\end{frame}

\begin{frame}{세 층이 부딪히는 지점 — 연결 규칙 셋}
\begin{itemize}
\item \textbf{스프린트 층은 '인수 FP'로, 월 층은 'EV'로 성과를 말한다} \ra{} 잇는 것은 Rules of Credit(인수 FP ÷ 620 × 1.96). 규칙이 없으면 리뷰의 "완료"와 보고의 "가치"가 다른 것을 가리킨다
\item \textbf{스프린트 층은 범위를 재조정하고 월 층은 범위를 동결한다} \ra{} 잇는 것은 접수 마감(동결 2주 전)과 변경협의체. 없으면 백로그 조정이 범위 변경의 뒷문
\item \textbf{스프린트 층의 지연은 다음 스프린트로 이월되지만 월 층의 지연은 임계경로 TF를 먹는다} \ra{} 잇는 것은 캘린더와 CPM이 같은 날짜를 가리키는 것
\item 그래서 S18 연장 결정(09-04)에 "임계경로 TF 영향" 칸이 있어야 했다 — 있었다면 예외 보고는 09-11에 나갔다(3주 빠름)
\item 조기 신호 2 — 5월 2주차 work item age 12영업일은 스프린트 층 지표가 월 층 대응(대체 결정권자)을 3주 당길 수 있었다는 기록
\end{itemize}
\keymsg{한 층만 있는 조직 — 스프린트 리뷰만 있어 임계경로가 안 보이거나, 월 보고만 있어 지연을 한 달 늦게 안다}
\note{§1.2 후반과 "실무에서 어떻게 틀리는가". 실무 오류 셋: 한 층만 있다 / 층 사이의 연결 규칙이 없다(스프린트 "완료"를 EV로 그대로 써서 착수 시 100\% 인정 — PV와 같아지게, 또는 백로그 조정을 변경 절차 밖에 둔다) / 아래층의 결정에 위층 영향 칸이 없다(스프린트 연장·항목 이월·시험 인력 이동은 일상 결정이지만 그 결정마다 임계경로 TF 영향이 계산되지 않으면 위층은 그 합을 예외 보고 시점에 처음 본다). 온담에서 스프린트 층의 "S10을 3주로" 결정이 월 층의 "A05 TF $-$5"를 만들고, 넉 달 뒤 S18의 $-$10과 합쳐져 게이트 층의 예외 보고가 되었다. 세 층은 각각 제 일을 했다 — 문제는 층 사이였다. 예상 소요 4분.}
\end{frame}

\begin{frame}{표준에서는 — PMBOK 8판 두 Focus Area · 다섯 성과영역 · ISO 21502 §7}
\begin{tbl}[\footnotesize]{L{0.2\textwidth}L{0.5\textwidth}L{0.2\textwidth}}
\textbf{8판 성과영역} & \textbf{실행·통제에서 무엇을 다루는가} & \textbf{이 교재} \\ \midrule
Governance & 변경 통제의 의사결정 구조(CCB), 보고·에스컬레이션, 예외 — 승인된 변경이 기준선과 PMB를 갱신 & 제3장 \\
Schedule & 일정 통제(CPM 갱신·TF 소진·마일스톤 전망) + 적응형 지표(번다운·처리량) 병기 & 2.8절 \\
Finance & 원가 통제와 EVM(EV·PV·AC \til{} EAC·ETC·TCPI·VAC), 예측, Earned Schedule 참고 & 제2장 \\
Risk & 트리거 관측 \ra{} 대응 실행 \ra{} 잔여·2차 평가 \ra{} 예비비 대조 & 4.5절 \\
Stakeholders & 이슈 관리와 에스컬레이션, 참여 감시 & 4.4·5.3절 \\
\end{tbl}
\begin{itemize}
\item 8판은 Quality 성과영역에서 QA·QC를 다시 구분하되 \textbf{QC의 결과가 인수의 근거}라는 연결을 앞세운다 — 온담의 "감리 지적은 검수의 근거"(회사 정본 §6.3)
\item ISO 21502 §7 — 변경 통제(7.10)·리스크(7.8)·\textbf{이슈(7.9, 별도 조항)}·보고(7.12)·조달(7.17). 계약이 중립적으로 인용할 수 있는 표준 [조항 번호는 원문 대조]
\end{itemize}
\keymsg{6판 대조(4.5 / 4.6 / 6.6 / 7.4 / 8.3 / 11.6 / 12.3 / 13.4)는 워크북 각 X.2 표에 — 사내 표준이 6판이면 6판 번호로 인용하고 8판 명칭을 병기}
\note{§1.3 PMBOK 8판·ISO 부분. 8판(2025)은 7판의 원칙·성과영역 위에 프로세스 기반 Focus Area를 되돌려 놓았고, 실행·통제에 해당하는 것은 Executing과 Monitoring and Controlling 두 Focus Area다. 6판의 8.3 Control Quality는 "QC = 산출물 검사, QA = 프로세스 검사"의 이분법이었고 7판이 그것을 지웠으며 8판이 다시 구분하되 인수와의 연결을 앞세운 것 — 온담에서 그 연결은 "심각도 1·2 미해결 0은 G3 인수 조건"이다(제4장). ISO 21502가 이슈 관리(7.9)를 리스크 관리(7.8)와 별도 조항으로 둔 것이 이 교재가 이슈 로그와 리스크 등록부를 한 문서의 두 면으로 다루되 분리하는 근거(4.4절). 실무 오류 — 6판 ITTO를 8판이라고 부르는 것. 예상 소요 5분.}
\end{frame}

\begin{frame}{PRINCE2 7 — 세 프로세스와 예외에 의한 관리(management by exception)}
\fig[0.56]{fig_1_3_management_by_exception}
\figcap{그림 1-3. 허용범위는 아래로 내려오는 위임이고 예외 보고는 위로 올라가는 반대급부다. 세 보고서(Checkpoint·Highlight·Exception)는 주기와 수신자가 다르며, 합치면 예외가 월례에 묻힌다}
\keymsg{예외 보고는 권한의 반납이 아니다 — "이 밖의 결정은 당신의 것이니 대안을 들고 가겠다"는 위임의 반대급부다}
\note{§1.3 PRINCE2 부분. 세 프로세스 — Controlling a Stage(PM의 일상: 작업 패키지 승인, 이슈·리스크 접수, Highlight Report, 시정 조치, 에스컬레이션), Managing Product Delivery(팀 매니저 = 수행사 PM 이재현: 작업 패키지 수락·실행·Checkpoint Report·인도), Managing a Stage Boundary(게이트 앞: 다음 단계 계획, BC 갱신, End Stage Report, 예외 계획). 분리가 주는 것 둘 — 발주사 PM과 수행사 PM의 역할이 문서 수준에서 갈라진다(Checkpoint는 이재현 \ra{} P1 PM, Highlight는 P1 PM \ra{} 정미란), 게이트가 단계를 닫고 여는 결정 지점이라는 것. Progress practice의 관리 방식 — 각 계층은 위가 준 허용범위 안에서 자율 결정하고, 넘을 것으로 예상되는 순간 예외 보고. 허용범위 없는 조직은 모든 결정이 위로 올라가 PM이 없고, 예외 보고 없는 조직은 PM이 권한 밖 결정을 혼자 내려 스폰서가 없다. Exception Report 구조(원인·6개 성과 목표 영향·대안 3·권고·결정 시한)가 헌장 §12 양식이고 EX-01(10-02)이 첫 실물. 실무 오류 — 예외 보고를 "실패 보고"로 이해해 내지 않는 것(늦을수록 대안이 줄고, 대안 없는 예외 보고는 사후 통보), Highlight를 매달 전부 G로 내는 것, Checkpoint와 Highlight를 합쳐 수행사 작업 패키지 상태가 스폰서에게 그대로 가는 것. 예상 소요 6분.}
\end{frame}

\begin{frame}{한국 SI의 통제 사이클 — 네 주기가 세 조직에 걸쳐 있다 · 분야별 대조}
\begin{itemize}
\item \textbf{주간·월간 보고}(발주사 PMO) · \textbf{감리}(제3자 — 온담 G1·중간 2회·G3·G4, 4.2억) · \textbf{검수}(계약·구매·재경 — §14 15영업일). 셋을 잇는 계약 문장: \textbf{"감리 지적사항은 검수의 근거"}
\item 그래서 감리 지적 하나가 검수를 보류시키고 \ra{} 검수 보류가 수행사의 대기 비용 청구가 되며 \ra{} 대기 비용의 귀책이 감리 보고서의 \textbf{귀속 칸}에서 결정된다(4.3절·5.5절)
\item 한국 SI에 EVM이 드문 이유 — 월간 보고의 "공정률"은 시간 비율이거나 담당자의 느낌이고, 대부분 계약 총액과 지급 일정만 있어 PMB가 없다
\end{itemize}
\begin{tbl}[\scriptsize]{L{0.13\textwidth}L{0.2\textwidth}L{0.3\textwidth}L{0.27\textwidth}}
\textbf{분야} & \textbf{통제 단위} & \textbf{성과 측정} & \textbf{인수·검수} \\ \midrule
건설·EPC & 기성(progress payment) & 물리적 진도 규칙이 계약 문서, EVM은 계약 요건(EIA-748 계열) & ITP hold point, NCR, 상주 감리, SCL 클레임 절차 \\
SI(한국) & 릴리스·단계 검수 & 공정률(시간 비율) — EV로 오해되기 쉬움. FP 인수가 자연스러운 EV 단위 & 감리 지적 \ra{} 검수 \ra{} 대금, 검수 지연 = 대기 비용 \\
제조 NPD & 게이트(DR·PV·양산 이관) & 마일스톤 가중 EV, 양산 원가 전망 & DV·PV, 시제품 결함, ECN이 형상관리와 직결 \\
제약 임상 & 실사·규제 제출 & 등록 환자·사이트 개시 수, 원가보다 일정 & GCP 실사, 규제 질의 — 대안이 "기다린다"가 되기 쉬움 \\
\end{tbl}
\keymsg{네 분야가 공유하는 것 — 인수의 조건이 통제 지표를 결정한다. 지표는 "표준이 권하는 것"이 아니라 "돈이 움직이는 조건"에서 거꾸로 정해진다}
\note{§1.4. 그림 1-4(한국 SI 네 주기와 세 조직, 09-25 감리 통보 \ra{} 10-16 검수 통보 연쇄)는 교재에 있으며 여기서는 표로 대신한다. 온담이 Day2에서 월별 PV 표를 만든 것은 "EVM을 하기 위해서"가 아니라 4월의 SPI 하락이 조달인지 개발인지를 열별로 구분하기 위해서였고(Day2 §4.6 해설), 그 예고가 실제로 일어났다. 실무 오류 — 공정률을 EV로 쓰는 것 / 감리와 검수를 다른 조직의 일로 두고 PM이 연결을 관리하지 않아 감리 보고서 귀속 칸이 검수 분쟁의 유일한 기록이 되는 것 / 주간 보고를 수행사가 쓰고 발주사가 읽기만 하는 것 — 발주사의 결정 SLA 준수율이 어느 문서에도 없으면 변경협의체에서 "발주사가 5월에 12일 미뤘다"에 반박할 자리가 없다(5.4절). 예상 소요 5분.}
\end{frame}

\begin{frame}{가장 흔하고 비싼 오류 셋 — 90\% 증후군 · 기준선 없는 보고 · 계획을 고쳐 편차를 없애기}
\fig[0.52]{fig_1_5_rebaseline_trap}
\figcap{그림 1-5. 왼쪽: BL-01 고정 + BL-02는 승인 변경(+1.02·+185 FP)만 계단으로 — A10 $-$10은 편차로 남는다. 오른쪽(가상): 분기 재기준선 — 마지막 기준선 대비 편차는 0이지만 최초 대비 누적 지연은 어느 문서에도 없다}
\keymsg{셋의 공통점 — 측정 규칙·기준선 버전·편차 이력이 문서에 없다. 통제는 분석 능력이 아니라 기록의 문제이고, 기록은 양식의 문제다}
\note{§1.5. 90\% 증후군 — 담당자 진척률은 종료 3개월 전까지 항상 80\til{}90\%. 거짓말이 아니라 측정 규칙의 부재("진척률"이 활동 수·기간·원가·FP·시험 통과율 중 무엇의 비율인지 정의되지 않음). Kerzner의 진척 착시(6.8절). 처방은 사람이 아니라 Rules of Credit(2.2절) — 진행 중 0(0/100), 물리 단위(인수 FP), 마일스톤 가중. 온담 SI 구현 EV가 "진행 중 스프린트 0"인 이유이며 그것이 만드는 구조적 SV $-$1.96은 착시가 아니라 규칙의 비용. 기준선 없는 보고 — "지난달보다 좋아졌다"는 추세이지 편차가 아니다(방향만 말하고 도착지를 말하지 않음); 버전 미명시(BL-01인가 02인가)는 EV와 PV가 다른 범위를 잴 수 있으니 제9호는 첫 표에 "비교 기준선 BL-01, BL-02 10-23 예정, 미반영 변경 CR-03·07 별도"를 적었다. 계획을 고쳐 편차 없애기 — 가장 비싼 오류. 재기준선이라 부르지만 편차의 은폐. "17개월 계획이 20개월 걸렸는데 마지막 기준선 대비 편차 0" 보고서가 이렇게 만들어진다. 재기준선의 조건(3.4절 — CCB + 이사회, 최대 2회, 편차 이력 보존). 생각해 볼 질문 5를 던진다 — 재기준선 0회 프로젝트와 3회 프로젝트 중 어느 쪽이 잘 관리된 것인가, 답하려면 무엇을 더 알아야 하는가. 예상 소요 5분.}
\end{frame}

\begin{frame}{EVM의 세 숫자와 그 앞의 하나 — PMB. ANSI/PMI 19-006-2019가 정의하는 것}
\begin{tbl}[\footnotesize]{L{0.2\textwidth}L{0.36\textwidth}L{0.36\textwidth}}
\textbf{표준 용어} & \textbf{정의} & \textbf{온담 P1} \\ \midrule
Control account & 범위·일정·원가가 통합되어 성과를 측정하는 단위 & WBS L2 다섯 산출물 + 원가 항목(SI·상용SW·인프라·단말·전환·시험·감리·PMO) \\
Work / Planning package & 측정 단위 / 아직 상세화되지 않은 원거리 예산 묶음 & 작업패키지 43개 / R4·R5 구간(rolling wave) \\
\textbf{PMB · BAC} & 시간대별로 배분된 예산의 합 — 관리예비비 제외 & \textbf{월별 PV 표 17개월, 합 145.80} \\
Undistributed budget & 승인되었으나 통제 계정에 미배분 & 승인 변경 CR-06·08·09의 1.005 — BL-02 편입 전까지 \\
Management reserve & PMB 밖, 미식별 리스크용, PM 권한 밖 & 관리예비비 12.0(정미란) \\
(Contingency) & 표준은 위치를 조직 관행에 맡긴다 & PMB 밖 · 집행 한도 안 · \textbf{집행 시 PV 없는 AC} \\
\end{tbl}
\keymsg{PV = 하기로 한 일의 예산 / EV = 한 일의 예산(계획 당시 값) / AC = 쓴 돈 — 셋이 같은 단위(돈)여서 일정과 원가가 한 표에 놓인다}
\note{§2.1 개념. "60\%의 기간이 지났는데 40\%의 가치가 인도되었다"는 문장은 시간과 가치가 같은 단위일 때만 성립한다 — EVM의 유일하고 결정적인 장점. 세 숫자 앞에 PMB가 있어야 한다 — Day2 4.3절이 월별 PV 표를 만든 이유. PMB 없으면 PV 없고, PV 없으면 SV·SPI 없다. 2019년 ANSI 표준은 2011년 Practice Standard 2판을 대체하고 적응형·증분 인도 개념을 통합한 현행 정본이다. 실무 오류 — BAC를 승인 예산 168.0으로 두어 CPI·TCPI가 좋아 보이는 것 / PMB 없이 계약 총액을 기간으로 균등 배분해 PV로 쓰는 것(4월 조달 봉우리 같은 구조가 사라져 편차 원인을 열별로 분리할 수 없다) / 관리예비비를 PMB에 넣는 것 / 컨틴전시 위치를 정하지 않아 집행이 어느 열의 AC로 나타나는지 매번 다른 것. 예상 소요 5분.}
\end{frame}

\begin{frame}{BAC는 145.8인가 156.0인가 — 온담이 실제로 부딪힌 질문}
\fig[0.56]{fig_2_1_pmb_bac_layers}
\figcap{그림 2-1. 컨틴전시 10.2는 PV 곡선 위 어디에도 배분되어 있지 않다. 온담은 BAC = 145.80(PMB), 예외 판정 기준 = 156.00(집행 한도)으로 분리했고 컨틴전시 집행은 PV 없는 AC로 기록한다 — CR-10이 헌장 문언으로 확정}
\keymsg{배분되지 않은 돈을 BAC에 넣으면 EV는 결코 BAC에 도달하지 못하고, SPI는 종료 시점에도 1이 되지 않으며, TCPI는 항상 실제보다 낮게 나온다}
\note{§2.1 "BAC는 145.8인가 156.0인가" — 워크북 ⑫ 해설 첫 항목. Day2 4.2절이 PMBOK 용어로 "원가 기준선 = 계획 원가 + 컨틴전시 = 156.0"이라 한 것은 PM이 책임지는 총액이라는 뜻이고 맞다. 그러나 EVM의 BAC는 시간대별로 배분된 PMB의 합이어야 하고 Day2 PV 표의 합은 145.80이다. 표준 용어로는 PMB(145.8)와 "관리예비비를 제외한 프로젝트 예산"(156.0)의 구분. 컨틴전시의 위치는 조직이 정한다 — PMB 안에 넣어 통제 계정별로 배분하는 조직도 있다. 온담은 밖에 두고 집행 시 PV 없는 AC로 기록하기로 했으므로 컨틴전시 집행 1.41 중 AC 발생분 1.19가 SI 열의 CV $-$1.19로 나타난다 — 원가 초과가 아니라 규칙의 결과이고, 그래서 SI 열의 원가 성과는 CPI가 아니라 "컨틴전시 소진율"로 읽는다(2교시 착시 2). CR-10이 이것을 헌장 §12 문언("초과 예상 판정 기준 = EAC가 집행 한도 156.0을 초과할 전망")으로 확정한다 — FP 0의 변경이 성과보고의 판정 기준을 바꾼 사례(3교시). 예상 소요 5분.}
\end{frame}

\begin{frame}{Rules of Credit — 지표는 측정이 아니라 합의의 산물이다}
\fig[0.58]{fig_2_2_rules_of_credit}
\figcap{그림 2-2. 같은 작업, 여섯 규칙, 여섯 EV 곡선. 0/100·마일스톤 가중·물리적 \%·LOE가 온담 채택 규칙, 50/50과 배분 노력은 미사용(라이선스의 50/50은 AC 인식 규칙이지 EV 규칙이 아니다). 아래 두 상자 = 규칙의 어긋남이 만드는 구조적 편차(1.96 · 4.70)}
\keymsg{SPI 0.943은 자연에 있는 값이 아니라 G1에서 합의한 규칙이 만든 값이다 — 규칙이 달랐으면 다른 값이다}
\note{§2.2. 그림은 위 여섯 곡선을 왼쪽부터 읽고, 아래 두 상자로 내려온다. 여섯 기법과 온담 적용 — 0/100(설계 스프린트 리뷰 인수 1.25/회, 리허설 회차, 감리·시험 회차 보고서) / 50/50(미사용) / 마일스톤 가중(단말 입고 검수 70 · 설치 확인 30, 통합·시험 스프린트 = 통합시험 TC 통과율 60 · UAT 종료 40) / 물리적 \%(구현 스프린트 = 인수 FP ÷ 620 × 1.96, 전환 매핑 = 확정 테이블 수 ÷ 1,842) / 배분 노력(미사용) / LOE(PMO 월 0.56, 유지보수, 클라우드 — EV = PV). 기법의 선택이 SPI를 가른다 — 느슨한 규칙(착수 시 100\%, 담당자 진척률)은 SPI가 1 근처에 붙어 있다가 종료 직전 무너지고, 엄격한 규칙(진행 중 0)은 구조적으로 1 아래에서 시작해 실제 편차를 더한 값으로 움직인다. 두 가지 결론 — 규칙은 계획 시점에 확정하고 실행 중 바꾸지 않는다(바꾸면 EV는 원하는 숫자가 된다) / 문서로 남긴다(원가 기준선 v1.0 부속 2 [추가 설정], BL-02에서 정식 첨부 — 워크북 §6 갱신 6번). 실무 오류 — SPI가 나쁠 때 "진행 중 50\% 인정"으로 완화하는 것, LOE 항목을 물리적 \%로 재려는 것(LOE는 SV = 0이므로 AC로만 읽는다), 규칙을 남기지 않아 감리가 "EV 산정 근거"를 물을 때 답이 없는 것. 예상 소요 7분.}
\end{frame}

\begin{frame}{PV 규칙과 EV 규칙의 시차 — 1.96 · EV 규칙과 AC 인식 기준의 시차 — 4.70}
\begin{itemize}
\item Day2 PV 배분 규칙 "스프린트 \textbf{시작 월}에 계상" vs EV 규칙 "스프린트 리뷰에서 \textbf{인수}된 FP를 인수 월에" — 각각 합리적(PV는 자금 조달 시점, EV는 인수 시점)
\item 둘이 어긋나면 월말에 걸친 스프린트 하나만큼 구조적 SV 음수 — 2026-09-30에 S20(09-28 착수)의 PV \textbf{1.96}. 정상 진행에서도 EV 0 \ra{} SV $-$4.40 중 $-$1.96은 지연이 아니다
\item 워크북 ⑫ 항목 4는 이것을 "조정 1: 측정 시차 제거 — SV $-$2.44, SPI 0.968"로 \textbf{별도 행}에 적는다. 병기이지 대체가 아니다
\item 같은 구조가 라이선스에 — EV 규칙 "인도·설치 확인서 100\%" vs AC 인식 "인보이스 발생액", 03-13 계약 변경으로 지급 조건이 인도 50\% / 통합시험 통과 50\% \ra{} EV 100\%, AC 50\% = CV \textbf{+4.70}
\item 이것은 EV 규칙의 문제가 아니라 \textbf{AC 인식 기준이 실행 중 바뀌었는데 양식에 조정 행이 없었던} 문제다
\end{itemize}
\keymsg{조정값을 본값 대신 쓰면 기준선 규칙을 실행 중에 바꾸는 것이다 — 옆에 적고, 이름을 붙이고, 크기를 적는다}
\note{§2.2 "온담의 사례로 읽기". 4교시 손계산에서 학생이 가장 자주 틀리는 자리 — 조정 SV를 구해 놓고 본값을 지운다. 미발생 4.70 = MDM 2차·APM 6.0 × 50\% + SAP 애드온 3.4 × 50\%(인도는 각각 5월·7월), 2027-01 통합시험 통과 시 한꺼번에 발생 — 그래서 2027-01 자금 봉우리(단말·DR·시험 16.20)에 얹혀 R-17의 일부가 된다. 학생에게 묻기 — "이재현 PM이 '진행 중 스프린트 50\% 인정'을 요구하면 무엇을 근거로 거절하는가"(제2장 생각해 볼 질문 4): 근거는 G1에 확정된 Rules of Credit 문서이고, 거절하지 않으면 SPI는 1 근처로 올라가 아무것도 재지 않게 된다. 예상 소요 4분.}
\end{frame}

\begin{frame}{편차와 지수 — 읽기 규칙 셋, 그리고 열별 분해는 의무다}
\begin{itemize}
\item 온담 2026-09-30 누적: PV 77.83 / EV 73.43 / AC 71.13 \ra{} SV $-$4.40 · CV +2.30 · SPI 0.943 · CPI 1.032
\item \textbf{SPI를 일수로 읽지 않는다} — 0.943은 비율. SV ÷ 월평균 PV 환산은 PV에 봉우리가 있으면 틀린다(4월 17.69 vs 9월 6.09, 세 배). 일수는 ES(2.4절) 또는 CPM(2.8절)이 준다
\item \textbf{월별 지수와 누적 지수를 섞지 않는다} — 누적 SPI 0.592 \ra{} 0.796 \ra{} 0.898 \ra{} \textbf{0.784(4월)} \ra{} 0.871 \ra{} 0.881 \ra{} 0.914 \ra{} 0.913 \ra{} 0.943. 회복의 일부는 수렴 착시
\item \textbf{열별로 분해한다}:
\end{itemize}
\begin{tbl}[\scriptsize]{L{0.08\textwidth}rrrrrrr}
 & \textbf{SI} & \textbf{상용SW} & \textbf{인프라} & \textbf{전환} & \textbf{시험·감리} & \textbf{PMO} & \textbf{합} \\ \midrule
SV & $-$5.31 & 0 & 0 & +0.20 & +0.70 & 0 & $-$4.40 \\
CV & $-$1.19 & \textbf{+4.70} & $-$1.00 & +0.20 & $-$0.05 & $-$0.36 & +2.30 \\
\end{tbl}
\keymsg{합계는 "일정은 조금 늦고 원가는 남는다" — 열별은 "개발이 늦고 감리 선행이 가리며, 원가는 라이선스 이연이 인프라·PMO 초과를 덮는다"}
\note{§2.3 "편차와 지수". 4월 0.784는 상용SW $-$6.0(조달 이연)이 SV $-$7.97의 75\% — 균등 배분 곡선이었다면 4월 SPI는 0.9쯤이었을 것이고 "약간의 지연"으로 넘어가 5월 R-05와 합쳐져 "지난달부터 조금씩 늦었다"는 잘못된 서사가 되었을 것이다. 봉우리 덕분에 4월은 조달(IS-01), 5월은 설계(IS-02)라는 두 사건으로 분리되었고 각각 다른 대응(지급 조건 협상 / 대체 결정권자)을 갖는다 — Day2 §4.6 해설의 예고가 실현된 자리. 표의 SI CV $-$1.19는 컨틴전시 집행 4건(R-07 0.15 + R-05 0.42 + R-03 0.50 + CR-03 0.12)이고, 인프라 $-$1.00은 클라우드 실사용(월 0.4 \ra{} 0.5, IS-10), PMO $-$0.36은 인건비·도구. 예상 소요 5분.}
\end{frame}

\begin{frame}{1교시 핵심 정리}
\begin{enumerate}
\item 통제 = 측정(같은 규칙) \ra{} 분석(원인) \ra{} 예측(전제) \ra{} 조치(세 층 — 격자가 가른다) \ra{} 기준선 갱신(승인된 변경만). 실무적 정체는 보고서 양식이다
\item 하이브리드 통제는 스프린트 / 월 / 게이트 세 층이고, 연결 규칙(Rules of Credit · 접수 마감 · 캘린더 = CPM 날짜)이 없으면 아래층의 작은 결정이 위층의 큰 사건이 될 때까지 안 보인다
\item PMBOK 8판 두 Focus Area · 다섯 성과영역, PRINCE2 7 세 프로세스 · 예외에 의한 관리(예외 보고는 위임의 반대급부), ISO 21502 §7(이슈와 리스크 별도), 한국 SI 네 주기 · "감리 지적은 검수의 근거"
\item 비싼 오류 셋 — 90\% 증후군 · 기준선 없는 보고 · 재기준선 남용 — 의 공통점은 기록의 부재
\item BAC = PMB의 합 145.80, 예외 기준 = 집행 한도 156.00, 컨틴전시는 PMB 밖·집행 시 PV 없는 AC. EV는 합의의 산물(Rules of Credit) — 구조적 시차 1.96·4.70은 병기하되 대체하지 않는다. SPI는 일수가 아니고, 열별 분해는 의무
\end{enumerate}
\keymsg{2교시는 "이대로 가면 어디에 도착하는가" — EAC 4공식, Earned Schedule, 다섯 착시, 그리고 임계경로 $-$10}
\note{교재 제1장 핵심 정리 1\til{}5와 제2장 핵심 정리 1\til{}2, 3의 전반. 생각해 볼 질문(제1장 1 — 여러분 회사 월간 보고서의 "진척률"은 무엇의 비율인가, 보고서 안에 정의가 있는가 / 제1장 3 — 예외 보고를 권한의 반납으로 느끼는 PM과 위임의 반대급부로 느끼는 PM의 행동은 어떻게 다른가, 여러분 조직에서 예외 보고를 낸 PM은 어떤 대우를 받았는가) 중 하나를 쉬는 시간 전에 던진다. 예상 소요 3분.}
\end{frame}
